\documentclass[11pt]{amsart}
\usepackage{amsmath,amssymb,enumitem,booktabs}
\newcommand{\eps}{\varepsilon}
\newtheorem{theorem}{Theorem}[section]
\newtheorem{lemma}[theorem]{Lemma}
\newtheorem{proposition}[theorem]{Proposition}
\newtheorem{corollary}[theorem]{Corollary}
\newtheorem{remark}[theorem]{Remark}
\newcommand{\pos}[1]{{(#1)_+}}
\DeclareMathOperator{\Tr}{Tr}
\title[Region II for small odd cycles]{The Region-II scalar inequality for the small odd cycles $3\le m\le 13$\\ (version 2)}
\begin{document}
\maketitle

\section{Statement, notation and result status}\label{sec:intro}

This note extends the Region-II scalar inequality of
\texttt{paper\_region2\_v2.tex} (cited below as [R2]) from odd $m\ge15$ to all
odd $3\le m\le 13$, so that the per-cycle theorems for $C_5,\dots,C_{13}$ can
be removed from the main chain.  All notation is from [R2]:
$q\in(\tfrac13,\tfrac12)$, $p=1-q$, $\alpha\in(q,r(q)]$ the frontier
eigenvalue, $L=\sqrt{pq-\alpha^2}$, $d=\alpha-q$, $e=1-2\alpha$,
$f=\alpha-L$, $\kappa=d/e$, $x=\alpha/p$, $\tau=q/\alpha$, $y=L/p$,
$\ell=L/\alpha$, $u=1-x$, $\delta=\alpha-\tfrac13$,
\[
 k_m(\lambda)=\frac{p^{m-1}-\lambda^{m-1}}{p+\lambda},\quad
 A_m=2L^{m-2}+mk_m(\alpha),\quad
 B_m=2L^{m-2}+mk_m(L),\quad
 R_m=\alpha^m+L^m-pq^{m-1},
\]
and $C_m$, $\xi$, $\rho$, $\psi(\xi,\rho)$ as in [R2, Def.~huber].  We use the
$(e,\kappa)$ chart of [R2, Lem.~chart]: the admissible $(q,\alpha)$ are
exactly $0<e<\tfrac13$, $0<\kappa\le\kappa_{\max}(e)=\frac{1-e}{1+e}$,
$\kappa<\kappa_q(e)=\frac{1-3e}{6e}$, and on it
$L^2=\alpha e-d(d+e)$, $\xi=(1-e)^2\kappa/e$, $\delta=\frac{1-3e}{6}$.

\begin{theorem}\label{thm:main}
For every odd $m$ with $3\le m\le13$ and every admissible $(q,\alpha)$,
\begin{equation}\label{eq:target}
 R_m\ \le\ C_m\,\psi(\xi,\rho).
\end{equation}
Consequently (Remark~\ref{rem:muniform}), the odd-cycle bound
$t(C_m,W)\ge p^m-pq^{m-1}$ holds in Region II for these $m$.
\end{theorem}

\begin{remark}[$m$-uniformity of the reduction in {[R2]}]\label{rem:muniform}
The statement of [R2, Thm.~huber] carries the standing hypothesis
``odd $m\ge15$'' of that paper's domain $D$, but its proof --- and the proofs
of the whole reduction chain it rests on, namely
[R2, Prop.~shift, Lem.~no-frontier, Lem.~one-frontier,
Lem.~forced-variance, Prop.~master-defect, Prop.~psi-forms] --- is
$m$-pointwise: $m$ enters only through the quantities $A_m$, $B_m$, $R_m$
defined above, and the only structural inputs are
$0<A_m<B_m$ [R2, Rem.~knegative] and the frontier geometry $L<q<\alpha<p$,
which hold for every odd $m\ge3$ on the admissible domain.
($A_m<B_m$ because $k_m$ is strictly decreasing on $[0,p]$ and $L<\alpha$;
$A_m>0$ because $\alpha<p$ gives $k_m(\alpha)>0$.)
Hence the implication
``\eqref{eq:target} on the admissible domain $\Rightarrow$
$t(C_m,W)\ge p^m-pq^{m-1}$ in Region II'' is valid verbatim for every odd
$m\ge3$, and Theorem~\ref{thm:main} may legitimately invoke it for
$3\le m\le13$.
\end{remark}

\emph{Status of the proof.}  Steps proved analytically below: the certificate
relaxation (Lemma~\ref{lem:relax}), the positivity gate
(Lemma~\ref{lem:gate}), the small-$e$ range $e\le1/60$
(Proposition~\ref{prop:smalle}), the tail cut on the Tur\'an band
(Lemma~\ref{lem:tail}: analytic reduction to six one-line exact rational
inequalities, displayed in Table~\ref{tab:tailints}), and the reduction of
the Tur\'an band $\delta\le\Delta_m$ to two explicit polynomial
nonnegativity claims (Proposition~\ref{prop:band}), which are then certified
by exact Bernstein coefficient tables.  The remaining moderate range
($1/60\le e\le\tfrac13-2\Delta_m$) is covered by an exact-rational certified
subdivision (Proposition~\ref{prop:moderate}), in the same certificate
protocol as [R2]; replacing these two certificate steps by prose arguments
is left open.

Throughout, a terminating decimal denotes that exact rational number, and
every decimal comparison is rounded in the safe direction with the
underlying exact comparison stated; all displayed inequalities are
additionally machine-verified by \texttt{validate\_smallm\_v2.py}.

\section{The relaxed certificate and the positivity gate}

\begin{lemma}[Relaxed $\lambda=1$ certificate]\label{lem:relax}
On the admissible domain, with $w:=\sqrt{2\alpha}$,
\begin{equation}\label{eq:rho-id}
 \rho=\frac{A_m w}{4B_m f},
\end{equation}
and consequently
\begin{equation}\label{eq:Tm}
 C_m\psi(\xi,\rho)\ \ge\ \sqrt{2\alpha}\,B_mf\Bigl(d-\frac{e^2}{16\alpha^2(1+\rho)}\Bigr)
 \ \ge\ \sqrt{2\alpha}\,B_mfd-\frac{e^2B_m^2f^2}{4\alpha^2A_m}.
\end{equation}
Hence \eqref{eq:target} follows from
\begin{equation}\label{eq:Tm2}
 (T_m):\qquad 4\alpha^2A_m R_m\ \le\ 4\alpha^2 A_m\sqrt{2\alpha}\,B_mfd-e^2B_m^2f^2 .
\end{equation}
\end{lemma}

\begin{proof}
By [R2, Def.~huber], $\rho=\frac{A_m}{B_m}\frac{\sqrt\alpha}{2\sqrt2 f}$, and
$\frac{\sqrt\alpha}{2\sqrt2}=\frac{\sqrt{2\alpha}}4$, giving \eqref{eq:rho-id}.
The first inequality in \eqref{eq:Tm} is the dual certificate $\lambda=1$
[R2, eq.~cert-II].  For the second, $1+\rho\ge\rho$ and \eqref{eq:rho-id} give
\[
 \sqrt{2\alpha}\,B_mf\cdot\frac{e^2}{16\alpha^2(1+\rho)}
 \le \sqrt{2\alpha}\,B_mf\cdot\frac{e^2}{16\alpha^2}\cdot\frac{4B_mf}{A_m\sqrt{2\alpha}}
 =\frac{e^2B_m^2f^2}{4\alpha^2A_m}. \qedhere
\]
\end{proof}

\begin{lemma}[Positivity gate; all odd $m\ge3$]\label{lem:gate}
If $R_m>0$ then
\begin{equation}\label{eq:gate}
 (m-1)\kappa\ >\ x\bigl(1+\kappa-\ell^{m-2}\bigr).
\end{equation}
\end{lemma}

\begin{proof}
By [R2, eq.~defect-form], $R_m=\alpha^{m-1}(\alpha-p\tau^{m-1})+L^m$, so
$R_m>0$ means $\tau^{m-1}<x(1+\ell^m)$, since
$L^m/(p\alpha^{m-1})=x\ell^m$.  Bernoulli's inequality gives
$\tau^{m-1}=(1-d/\alpha)^{m-1}\ge1-(m-1)d/\alpha$, whence
\[
 (m-1)\frac d\alpha\ >\ 1-x-x\ell^m .
\]
Multiply by $\alpha/e$ and use $1-x=u=\frac{d+e}p=\frac{x(d+e)}\alpha$, so
that $(1-x)\alpha/e=x(1+\kappa)$; finally
$\ell^m\alpha/e\le\ell^{m-2}$ because $\ell^2\le e/\alpha$
[R2, Lem.~elementary].
\end{proof}

\begin{corollary}[Gate on the small-$e$ range]\label{cor:gate-smalle}
Let $3\le m\le13$ be odd, $e\le1/60$, and $R_m>0$.  Then, with the standing
bounds $x\ge\frac{1-e}{1+3e}\ge\frac{59}{63}>0.9365$ and
$\ell\le\sqrt{2e/(1-e)}\le1.4262\sqrt e\le0.1842$,
\begin{equation}\label{eq:gate-smalle}
 \kappa\ >\ \frac{x(1-\ell)}{m-1-x}\ \ge\ \frac{0.763}{m-1}\ \ge\ 0.0635,
 \qquad
 \xi=\frac{(1-e)^2\kappa}e\ \ge\ 3.68,
\end{equation}
and for $m=3$ moreover $\kappa>0.7183$.  In particular the region
$\{\,e\le1/60,\ \xi\le1,\ R_m>0\,\}$ is empty for $3\le m\le13$, and
\begin{equation}\label{eq:eps-smalle}
 \eps:=\frac{e^2}{16\alpha^2(1+\rho)d}=\frac{e}{4(1-e)^2(1+\rho)\kappa}
 \ \le\ \min\Bigl(0.068,\ \frac{e}{4(1-e)^2\kappa}\Bigr).
\end{equation}
\end{corollary}

\begin{proof}
Since $m\ge3$ and $\ell<1$, $\ell^{m-2}\le\ell$, so \eqref{eq:gate} gives
$\kappa(m-1-x)>x(1-\ell)$, i.e.\ the first bound; then
$x(1-\ell)\ge0.9365\cdot0.8158>0.763$ and $m-1-x<m-1$.  For $m=3$,
$t\mapsto t/(2-t)$ is increasing, so
$\kappa>0.9365\cdot0.8158/(2-0.9365)>0.7183$.  Next,
$\xi\ge(59/60)^2\cdot\frac{0.763}{12}\cdot60>3.68$ using $m\le13$; in
particular $\xi>1$.  Finally $\eps=1/(4\xi(1+\rho))\le1/(4\cdot3.68)<0.068$
and $\rho\ge0$; the equality in \eqref{eq:eps-smalle} is the chart identity
$\xi=(1-e)^2\kappa/e$.
\end{proof}
(The displayed decimal replacements: $0.9365\cdot0.8158=0.76399\ldots>0.763$,
$0.9365\cdot0.8158/1.0635=0.71837\ldots>0.7183$,
$(59/60)^2\cdot0.763\cdot5=3.68889\ldots>3.68$, and
$1/14.72=0.06793\ldots<0.068$.)

\section{The small-$e$ range $e\le1/60$}

\begin{proposition}\label{prop:smalle}
Let $3\le m\le13$ be odd, $e\le1/60$, $R_m>0$.  Then
$R_m\le\sqrt{2\alpha}B_mf\bigl(d-\frac{e^2}{16\alpha^2(1+\rho)}\bigr)\le C_m\psi$.
\end{proposition}

\begin{lemma}[Reduction; m-adapted {[R2, Lem.~reduceA]}]\label{lem:reduce}
Under the hypotheses of Proposition~\ref{prop:smalle}, the first inequality
follows from
\begin{equation}\label{eq:red}
 \sqrt{1-e}\,\frac{(1-\ell)(1-y^{m-1})(1-\eps)}{1+y}
 \ \ge\
 x^{m-2}\Bigl[\frac{m-1}{mx}-\frac{1+1/\kappa}{m}\Bigr]+\Lambda,
 \qquad
 \Lambda:=\frac{x^{m-2}(e/\alpha)^{m/2-1}}{m\kappa}.
\end{equation}
\end{lemma}

\begin{proof}
Exactly as in [R2, Lem.~reduceA]: $B_m\ge mk_m(L)=\frac{mp^{m-2}(1-y^{m-1})}{1+y}$,
$\sqrt{2\alpha}=\sqrt{1-e}$, $f=\alpha(1-\ell)$ give: payment
$\ge md\alpha p^{m-2}\cdot\mathrm{LHS}\eqref{eq:red}$; Bernoulli
$\tau^{m-1}\ge1-(m-1)d/\alpha$ and $L^m\le(\alpha e)^{m/2}$ give
$R_m\le md\alpha p^{m-2}\cdot\mathrm{RHS}\eqref{eq:red}$.
(No $m\ge15$ constant enters.)
\end{proof}

\begin{proof}[Proof of Proposition~\ref{prop:smalle}]
Standing bounds for $e\le1/60$, each stated with its exact justification:
$x\ge\frac{1-e}{1+3e}\ge\frac{59}{63}>0.9365$ (from $d\le e$);
$\ell^2\le e/\alpha=\frac{2e}{1-e}\le\frac{120}{59}e\le2.04e$, hence
$\ell\le1.4262\sqrt e\le0.1842$
($1.4262^2=2.03404\ldots>120/59=2.03389\ldots$;
$1.4262/\sqrt{60}<0.1842$); $y\le\ell$ (as $p\ge\alpha$);
$\sqrt{1-e}\ge1-0.51e$ (as $(1-0.51e)^2-(1-e)=-0.02e+0.2601e^2\le0$ for
$e\le0.02/0.2601$); $1/x=1+(1+\kappa)e/\alpha\le1+2.04(1+\kappa)e$
(as $1/\alpha\le120/59\le2.04$);
$u=1-x=\frac{(1+\kappa)e}{p}\ge1.9047(1+\kappa)e$
(as $p=\frac{1+e}2+d\le\frac{1+3e}2\le\frac{21}{40}$ and
$40/21=1.90476\ldots\ge1.9047$); and Corollary~\ref{cor:gate-smalle}.

\emph{Case $m=3$.}  Using
$\frac{(1-\ell)(1-y^2)(1-\eps)}{1+y}\ge(1-\ell)(1-y)(1-y^2)(1-\eps)$ and
$(1-a_1)\cdots(1-a_k)\ge1-\sum a_i$ on $[0,1]$,
\[
 \mathrm{LHS}\eqref{eq:red}\ \ge\ 1-2\ell-y^2-\eps-0.51e
 \ \ge\ 1-0.3684-0.034-0.068-0.0085\ >\ 0.52,
\]
where $y^2\le\ell^2\le2.04e\le2.04/60=0.034$.  On the other side, for $m=3$
the bracket in \eqref{eq:red} is $\frac{2}{3x}-\frac{1+1/\kappa}3$, so
$\mathrm{RHS}\eqref{eq:red}=\frac23-\frac{x(1+1/\kappa)}3+\Lambda$ with
$\Lambda=\frac{x\sqrt{e/\alpha}}{3\kappa}$.  Using $\kappa>0.7183$,
$1+1/\kappa>2$, $x\ge0.9365$, $x\le1$, $\sqrt{e/\alpha}\le1.4262\sqrt e$
and $1.4262/(3\cdot0.7183)=0.66183\ldots\le0.662$:
\[
 \mathrm{RHS}\eqref{eq:red}\le\frac23-\frac{0.9365\cdot2}{3}+0.662\sqrt e
 \le\frac23-0.6243+0.662\cdot0.1291<0.128<0.52,
\]
($\sqrt{1/60}<0.1291$ as $0.1291^2=0.01666\,681>1/60$).  So \eqref{eq:red}
holds for $m=3$.

\emph{Case $5\le m\le13$.}  Put $z:=(m-2)u$, $c_m:=\frac{m-3}{2(m-2)}$,
$\beta_m:=\frac{1.0711}{m-2}$ (so that $2.04(1+\kappa)e\le\beta_m z$, since
$\beta_mz\ge1.0711\cdot1.9047(1+\kappa)e=2.04012\ldots(1+\kappa)e$), and
\[
 \Gamma_0(e):=2.8524\sqrt e+0.51e+(2.04e)^2,\qquad
 \sigma_m(e):=1.21\,(2.04e)^{(m-2)/2}.
\]
Then $2\ell+0.51e+y^{m-1}\le\Gamma_0$ (using
$y^{m-1}\le\ell^4\le(2.04e)^2$ for $m\ge5$), and
$\Lambda\le\sigma_mx^{m-2}$: indeed
$(e/\alpha)^{(m-2)/2}\le(2.04e)^{(m-2)/2}$, and by
Corollary~\ref{cor:gate-smalle} and $m\le13$,
\[
 \frac1{m\kappa}\ <\ \frac{m-1}{0.763\,m}\ \le\ \frac{12}{0.763\cdot13}
 \ =\ 1.20979\ldots\ <\ 1.21 .
\]
Since $x^{m-2}\le(1+z+c_mz^2)^{-1}$ [R2, eq.~E-a] and
$1/x\le1+\beta_mz$, $\frac1{mx}\ge\frac1m$, \eqref{eq:red} follows from
\begin{equation}\label{eq:F}
 F:=\bigl(1-\Gamma_0-\eps\bigr)(1+z+c_mz^2)-\frac1x+\frac{2+1/\kappa}m-\sigma_m\ \ge\ 0 .
\end{equation}
We record two endpoint evaluations, rounded up
($2.8524/\sqrt{60}<2.8524/7.7459<0.36825$, since $7.7459^2=59.99896681<60$;
$0.51/60=0.0085$; $(2.04/60)^2=0.001156$; $1.293/60=0.02155$):
\begin{equation}\label{eq:G-ends}
 \Gamma_0(e)\le\Gamma_0(1/60)<0.378,\qquad
 \Gamma_1(e):=\Gamma_0(e)+1.293e\le\Gamma_1(1/60)<0.3995,
\end{equation}
both $\Gamma_0,\Gamma_1$ being increasing in $e$; and
$\sigma_m(e)\le\sigma_5(1/60)=1.21\cdot(0.034)^{3/2}<0.0076$ for all
$5\le m\le13$ (as $2.04e\le0.034<1$).

\emph{Sub-case $\kappa<\tfrac15$.}  Here $\frac{2+1/\kappa}m\ge\frac7m$, and
by \eqref{eq:eps-smalle} $\eps\le0.068$.  Since
$1-\Gamma_0-0.068-\beta_m\ge1-0.378-0.068-0.35704>0.19>0$
($\beta_m\le\beta_5=1.0711/3=0.35703\ldots$), dropping $c_mz^2\ge0$,
\[
 F\ \ge\ z\bigl(1-\Gamma_0-0.068-\beta_m\bigr)+\frac7m-\Gamma_0-0.068-\sigma_m
 \ \ge\ \frac7{13}-0.378-0.068-0.0076\ >\ 0.084\ >\ 0 .
\]

\emph{Sub-case $\kappa\ge\tfrac15$.}  Here \eqref{eq:eps-smalle} gives
$\eps\le\frac{5e}{4(1-e)^2}\le\frac{4500}{3481}e\le1.293e$, so
$1-\Gamma_0-\eps\ge1-\Gamma_1>0$ by \eqref{eq:G-ends}.  Dropping
$c_mz^2\ge0$ and using $1/x\le1+\beta_mz$, $\frac{2+1/\kappa}m\ge\frac3m$
($\kappa<1$ on the domain):
\[
 F\ \ge\ z\bigl(1-\Gamma_1-\beta_m\bigr)+\frac3m-\Gamma_1-\sigma_m .
\]
The bracket is positive ($\ge1-0.3995-0.35704>0.24$), and
$z\ge1.9047\cdot1.2\,(m-2)e\ge2.2856(m-2)e$ (from $1+\kappa\ge1.2$), so
\begin{equation}\label{eq:Gm}
 F\ \ge\ G_m(e):=2.2856(m-2)e\bigl(1-\Gamma_1(e)-\beta_m\bigr)
 +\frac3m-\Gamma_1(e)-\sigma_m(e) .
\end{equation}
Bounding the bracket once more by $\Gamma_1\le0.3995$ and expanding
$\Gamma_1(e)=2.8524\sqrt e+1.803e+4.1616e^2$,
\begin{equation}\label{eq:Ghat}
 G_m(e)\ \ge\ \widehat G_m(e):=\frac3m+\bigl(c_{1,m}e-2.8524\sqrt e\bigr)
 -4.1616e^2-\sigma_m(e),
\end{equation}
\begin{equation}\label{eq:c1}
 c_{1,m}:=2.2856(m-2)\bigl(0.6005-\beta_m\bigr)-1.803
 \ =\ 1.3725028\,(m-2)-4.25110616,
\end{equation}
an exact terminating decimal for each $m$ (Table~\ref{tab:smalle}).

\emph{$\widehat G_m$ is strictly decreasing on $(0,1/60]$.}  The three terms
$-4.1616e^2$, $-\sigma_m(e)$, and the constant are nonincreasing in $e$; for
the remaining term, its derivative satisfies, for $0<e\le1/60$,
\[
 \frac{d}{de}\bigl(c_{1,m}e-2.8524\sqrt e\bigr)
 =c_{1,m}-\frac{1.4262}{\sqrt e}
 \ \le\ c_{1,m}-1.4262\sqrt{60}
 \ <\ 10.847-11.045\ <\ 0,
\]
because $c_{1,m}$ is increasing in $m$ with
$c_{1,13}=10.84642464<10.847$, and
$1.4262\sqrt{60}>1.4262\cdot7.745=11.045919>11.045$
($7.745^2=59.985025<60$).  Hence
$\widehat G_m(e)\ge\widehat G_m(1/60)$ on $(0,1/60]$, and evaluating with
the safe roundings above ($2.8524\sqrt{1/60}<0.36825$,
$4.1616/3600<0.001156\,1$, $\sigma_m(1/60)\le\bar\sigma_m$ of
Table~\ref{tab:smalle}, $c_{1,m}\ge\underline c_{1,m}$ of
Table~\ref{tab:smalle}):
\[
 \widehat G_m(1/60)\ \ge\ \frac3m+\frac{\underline c_{1,m}}{60}
 -0.36825-0.0011561-\bar\sigma_m\ \ge\ \widehat G_m^{\,\flat}
\]
with the values $\widehat G_m^{\,\flat}$ of Table~\ref{tab:smalle}, all
positive (minimum $0.0383$ at $m=11$).  This proves $F\ge0$ in both
sub-cases, hence \eqref{eq:red}, hence the proposition via
Lemma~\ref{lem:reduce} and Lemma~\ref{lem:relax}.
\end{proof}

\begin{table}\centering
\begin{tabular}{lccccc}\toprule
$m$&5&7&9&11&13\\\midrule
$\beta_m$&$0.35704$&$0.21422$&$0.15302$&$0.11902$&$0.09738$\\
$\underline c_{1,m}$ &$-0.13360$&$2.61140$&$5.35641$&$8.10141$&$10.84642$\\
$\bar\sigma_m$&$0.0076$&$0.000258$&$0.0000088$&$0.0000003$&$0.00000002$\\
$\widehat G_m^{\,\flat}$&$0.2207$&$0.1024$&$0.0531$&$0.0383$&$0.0421$\\
\bottomrule
\end{tabular}
\caption{Small-$e$ constants.  $\beta_m$ is rounded up (safe in
$1-\Gamma_1-\beta_m$ and in $c_{1,m}$ via \eqref{eq:c1});
$\underline c_{1,m}$ is the exact value \eqref{eq:c1} truncated toward
$-\infty$ at the fifth decimal; $\bar\sigma_m$ is an upper bound for
$\sigma_m(1/60)=1.21(0.034)^{(m-2)/2}$; $\widehat G_m^{\,\flat}$ is the
resulting lower bound for $\widehat G_m(1/60)$, truncated down.}
\label{tab:smalle}
\end{table}

\section{The Tur\'an band $\delta\le\Delta_m$: polynomial Bernstein certificates}

Here $\delta=\frac{1-3e}6\in(0,\Delta_m]$ with
$\Delta_m=\tfrac1{40}$ for $m\le9$ and $\Delta_m=\tfrac1{80}$ for
$m\in\{11,13\}$; write $s:=1-d/\delta\in[0,1)$, so
$\alpha=\tfrac13+\delta$, $e=\tfrac13-2\delta$, $q=\tfrac13+s\delta$,
$p=\tfrac23-s\delta$, and
$L^2=\tfrac19-\tfrac{(2-s)\delta}3-(1+s^2)\delta^2$, a polynomial in
$(\delta,s)$.  The band covers all admissible $\kappa$ (there
$\kappa_q<\kappa_{\max}$).

\begin{lemma}[Chord]\label{lem:chord}
$\sqrt{2\alpha}\ \ge\ w_\ell(\delta):=0.8164+1.2\,\delta$ for
$0\le\delta\le\tfrac1{40}$.
\end{lemma}

\begin{proof}
$w_\ell(\delta)^2-(2/3+2\delta)$ is a quadratic in $\delta$ with positive
leading coefficient, hence maximal at the endpoints, where it is
negative: $0.8164^2=0.66650\ldots<2/3$ and
$(0.8164+0.03)^2=0.71639296<2/3+1/20=0.71666\ldots$.
\end{proof}

\begin{lemma}[Tail]\label{lem:tail}
Let $t_m^*$ be as in Table~\ref{tab:band}.  If $R_m>0$ and
$\delta\le\Delta_m$, then $d/\delta>t^*_m$, i.e.\ $s<1-t^*_m$.
\end{lemma}

\begin{proof}
\emph{Step 1 (gate $\Rightarrow$ curve).}
By Lemma~\ref{lem:gate}, $\kappa(m-1-x)>x(1-\ell^{m-2})$; on the band
$x=\alpha/p\ge\tfrac12$ (as $\alpha\ge\tfrac13\ge p/2$) and $m-1-x<m-1$,
and $\ell^{m-2}\le w(\delta)^{\nu}$ where
\[
 w(\delta):=\frac e\alpha=\frac{1-6\delta}{1+3\delta},\qquad
 \nu:=\frac{m-2}2
\]
(from $\ell^2\le e/\alpha$, [R2, Lem.~elementary]).  Hence, for $R_m>0$,
\begin{equation}\label{eq:phi}
 \frac d\delta=\frac{\kappa e}\delta
 \ >\ \phi_m(\delta):=\frac{e}{2(m-1)\delta}\,h(\delta),
 \qquad h(\delta):=1-w(\delta)^{\nu}.
\end{equation}

\emph{Step 2 (endpoint reduction: $\phi_m(\delta)\ge\phi_m(\Delta_m)$ on
$(0,\Delta_m]$).}
We claim $h$ is concave on $[0,\Delta_m]$.  Writing
$w=-2+3/(1+3\delta)$, so $w'=-9/(1+3\delta)^2$ and $w''=54/(1+3\delta)^3$,
\[
 (w^{\nu})''=\nu w^{\nu-2}\bigl[(\nu-1)w'^2+w\,w''\bigr]
 =\nu w^{\nu-2}\,\frac{27\,\bigl(3\nu-1-12\delta\bigr)}{(1+3\delta)^4}
 \ \ge\ 0
\]
on $[0,\Delta_m]$: indeed $w>0$ there (as $\delta\le\tfrac1{40}<\tfrac16$),
and $3\nu-1-12\delta\ge0$ because for $m\ge5$, $3\nu\ge\tfrac92>1+\tfrac{12}{40}$,
while for $m=3$, $3\nu=\tfrac32\ge1+12\delta$ iff $\delta\le\tfrac1{24}$,
and $\Delta_3=\tfrac1{40}<\tfrac1{24}$.  So $w^{\nu}$ is convex, $h$ is
concave, and $h(0)=0$; concavity then gives, for $0<\delta\le\Delta_m$,
\[
 h(\delta)\ \ge\ \tfrac{\delta}{\Delta_m}h(\Delta_m)
 +\bigl(1-\tfrac{\delta}{\Delta_m}\bigr)h(0)
 =\tfrac{\delta}{\Delta_m}h(\Delta_m),
 \qquad\text{i.e.}\qquad
 \frac{h(\delta)}\delta\ \ge\ \frac{h(\Delta_m)}{\Delta_m}.
\]
Combining with $e=\tfrac13-2\delta\ge\tfrac13-2\Delta_m$ (both factors
positive),
\[
 \phi_m(\delta)=\frac{e}{2(m-1)}\cdot\frac{h(\delta)}{\delta}
 \ \ge\ \frac{\tfrac13-2\Delta_m}{2(m-1)}\cdot\frac{h(\Delta_m)}{\Delta_m}
 \ =\ \phi_m(\Delta_m).
\]

\emph{Step 3 (endpoint comparison, exact).}
It remains to check $\phi_m(\Delta_m)>t^*_m$, i.e.\ with the rationals
\[
 c_m:=\frac{\tfrac13-2\Delta_m}{2(m-1)\Delta_m},\qquad
 w_m:=w(\Delta_m),\qquad r_m:=1-\frac{t^*_m}{c_m},
\]
that $w_m^{\nu}<r_m$.  Table~\ref{tab:tailints} lists $c_m$, $w_m$, $r_m$;
each $r_m\in(0,1)$, so squaring gives the equivalent \emph{rational}
inequality $w_m^{m-2}<r_m^2$, i.e.\ the one-line integer comparison in the
last column of Table~\ref{tab:tailints}, which is verified exactly (and can
be checked by hand or by any exact integer arithmetic).  This proves
$d/\delta>\phi_m(\delta)\ge\phi_m(\Delta_m)>t^*_m$.
\end{proof}

\begin{table}\centering
\small
\begin{tabular}{cccccl}\toprule
$m$&$\Delta_m$&$w_m$&$c_m$&$r_m$&integer form of $w_m^{m-2}<r_m^2$\\\midrule
$3$&$1/40$&$34/43$&$17/6$&$757/850$&$34\cdot850^2=24\,565\,000<24\,641\,107=43\cdot757^2$\\
$5$&$1/40$&$34/43$&$17/12$&$299/425$&$34^3\!\cdot\!425^2=7\,099\,285\,000<7\,108\,005\,307=43^3\!\cdot\!299^2$\\
$7$&$1/40$&$34/43$&$17/18$&$481/850$&$34^5\!\cdot\!850^2=32\,827\,093\,840\,000<34\,012\,020\,380\,923=43^5\!\cdot\!481^2$\\
$9$&$1/40$&$34/43$&$17/24$&$191/425$&$34^7\!\cdot\!425^2=9\,487\,030\,119\,760\,000<9\,916\,214\,751\,794\,467=43^7\!\cdot\!191^2$\\
$11$&$1/80$&$74/83$&$37/30$&$223/370$&$74^9\!\cdot\!370^2<83^9\!\cdot\!223^2$ \ (see below)\\
$13$&$1/80$&$74/83$&$37/36$&$493/925$&$74^{11}\!\cdot\!925^2<83^{11}\!\cdot\!493^2$ \ (see below)\\
\bottomrule
\end{tabular}
\caption{Tail-lemma endpoint data.  All entries are exact rationals:
$w_m=e(\Delta_m)/\alpha(\Delta_m)$, $c_m$ as in Step~3, $r_m=1-t^*_m/c_m$.
The two large comparisons are
$74^9\cdot370^2=9\,109\,382\,235\,108\,373\,145\,600
<9\,296\,351\,954\,199\,516\,700\,787=83^9\cdot223^2$ and
$74^{11}\cdot925^2=311\,768\,606\,996\,584\,070\,908\,160\,000
<313\,006\,138\,444\,263\,224\,418\,858\,083=83^{11}\cdot493^2$.}
\label{tab:tailints}
\end{table}

\begin{proposition}[Band]\label{prop:band}
For odd $3\le m\le13$, $R_m>0$ and $\delta\le\Delta_m$ imply $(T_m)$.
\end{proposition}

\begin{proof}
Write $A_m=A^e+LA^o$, $B_m=B^e+LB^o$, $R_m=R^e+LR^o$, $f=\alpha-L$ with
$A^e,\dots$ polynomials in $(\delta,s)$ (possible since $m-1$ is even and
$p^{m-1}-L^{m-1}$ is divisible by $p^2-L^2$).  Substituting
$w_\ell\le\sqrt{2\alpha}$ (Lemma~\ref{lem:chord}) into \eqref{eq:Tm2} and
collecting $L$-parities, $(T_m)$ follows from $G:=P+LQ\ge0$, where
\[
 P+LQ\ :=\ w_\ell(\delta)\,4\alpha^2d\,(A B f)\big|_{\text{parity}}
 -e^2\,(B f)^2\big|_{\text{parity}}-4\alpha^2 (A R)\big|_{\text{parity}},
\]
two explicit polynomials in $(\delta,s)$ with rational coefficients
(generated and archived by the validation script).  Let
$H:=L^2Q^2-P^2$; then $H=c_m\,\delta\,K$ with $c_m>0$ and $K$ a polynomial.
On the box $\mathcal B_m:=[0,\Delta_m]\times[0,1-t^*_m]$ (which by
Lemma~\ref{lem:tail} contains all band points with $R_m>0$):
\emph{all two-dimensional Bernstein coefficients of $Q$ and of $K$ on
$\mathcal B_m$ are nonnegative} (exact rational computation; no subdivision
is needed).  Hence $Q\ge0$ and $H\ge0$ on $\mathcal B_m$; if $P\ge0$ then
$G\ge0$ trivially, and if $P<0$ then $L^2Q^2\ge P^2$ gives $LQ\ge|P|$, so
$G\ge0$.  By \eqref{eq:Tm2} and Lemma~\ref{lem:relax} this proves the claim.
\end{proof}

\begin{table}\centering
\begin{tabular}{lcccccc}\toprule
$m$&3&5&7&9&11&13\\\midrule
$\Delta_m$&$1/40$&$1/40$&$1/40$&$1/40$&$1/80$&$1/80$\\
$t^*_m$&$0.31$&$0.42$&$0.41$&$0.39$&$0.49$&$0.48$\\\bottomrule
\end{tabular}\caption{Band parameters.}\label{tab:band}
\end{table}

\section{The moderate range: certified subdivision}

\begin{proposition}\label{prop:moderate}
For odd $3\le m\le13$ and $e\in[\tfrac1{60},\tfrac13-2\Delta_m]$, either
$R_m\le0$ or $(T_m)$ holds.
\end{proposition}

\begin{proof}[Certificate]
Exact-rational interval subdivision over
$(e,\kappa)\in[\tfrac1{60},\tfrac13-2\Delta_m]\times[0,1]$: on each box,
directed rational bounds for $\alpha,d,q,p,L^2,f,k_m(\alpha),k_m(L),A_m,B_m,
R_m$ and integer-square-root brackets for $L$, $\sqrt{2\alpha}$ certify
either $R_m\le0$ or \eqref{eq:Tm2}; boxes wholly outside the admissible
domain ($q\le\tfrac13$ or $\kappa>\kappa_{\max}$) are discarded.  The run
(same protocol as [R2, \S certificate-protocol]) terminates with
$6757,3911,3567,4307,6487,8785$ boxes for $m=3,5,\dots,13$.  This step is
\emph{computational}; a prose replacement is open.
\end{proof}

\section{Assembly}

\begin{proof}[Proof of Theorem~\ref{thm:main}]
If $R_m\le0$, \eqref{eq:target} is trivial ($\psi\ge0$).  If $R_m>0$:
for $e\le1/60$, Proposition~\ref{prop:smalle} applies (Corollary
\ref{cor:gate-smalle} shows this range meets $\xi>1$);
for $1/60\le e\le\tfrac13-2\Delta_m$, Proposition~\ref{prop:moderate}
gives $(T_m)$; for $e\ge\tfrac13-2\Delta_m$ (i.e.\ $\delta\le\Delta_m$),
Proposition~\ref{prop:band} gives $(T_m)$.  In the latter two cases
Lemma~\ref{lem:relax} upgrades $(T_m)$ to \eqref{eq:target}.
Finally, Remark~\ref{rem:muniform} converts \eqref{eq:target} into the
odd-cycle bound in Region II.
\end{proof}

\begin{remark}[Computational components]\label{rem:comp}
Two steps are certificate-based rather than prose: (i) the Bernstein
coefficient tables of Proposition~\ref{prop:band} (a classical positivity
certificate: nonnegativity of all Bernstein coefficients of an explicit
rational polynomial on a box implies nonnegativity of the polynomial), and
(ii) the subdivision of Proposition~\ref{prop:moderate}.  Both are
exact-rational and small.  Everything else is analytic; in particular the
tail cut (Lemma~\ref{lem:tail}) is analytic up to the six displayed
one-line integer comparisons of Table~\ref{tab:tailints}.  The
per-cycle theorems $C_5,\dots,C_{13}$ are no longer used.
\end{remark}

\end{document}
